\documentclass[11pt]{article}
\usepackage[utf8]{inputenc}
\usepackage{amsmath, amssymb, amsthm, geometry}
\usepackage{algorithm}
\usepackage{algorithmic}
\geometry{letterpaper, margin=1in}

\title{Analysis of Stochastic Gradient Descent for Smooth Non-Convex Functions}
\author{}
\date{}

\begin{document}
\maketitle

\section*{Algorithm Name}
Stochastic Gradient Descent (SGD)

\section*{Mathematical Formulation}
The algorithm iteratively updates the parameter vector $w \in \mathbb{R}^d$ by moving in the direction of the negative stochastic gradient. While originally described as using a subgradient for convex and Lipschitz functions, under the specified assumptions of smoothness and non-convexity, the subgradient coincides with the gradient.

Let $f: \mathbb{R}^d \rightarrow \mathbb{R}$ be the objective function defined as an expectation over a data distribution $\mathcal{D}$:
$$f(w) = \mathbb{E}_{z \sim \mathcal{D}}[f(w; z)]$$

The update rule at iteration $t$ is given by:
$$w_{t+1} = w_t - \eta_t g_t$$
where:
\begin{itemize}
    \item $w_t$ is the parameter vector at step $t$.
    \item $\eta_t > 0$ is the learning rate (step size) at step $t$.
    \item $z_t \sim \mathcal{D}$ is a single randomly sampled data point at step $t$.
    \item $g_t = \nabla f(w_t; z_t)$ is the stochastic gradient computed on $z_t$.
\end{itemize}

\section*{Pseudocode}
\begin{algorithm}[H]
\caption{Stochastic Gradient Descent (SGD)}
\begin{algorithmic}[1]
\REQUIRE Initial parameters $w_0 \in \mathbb{R}^d$, learning rate schedule $\{\eta_t\}_{t=0}^{T-1}$, number of iterations $T$
\FOR{$t = 0, 1, \dots, T-1$}
    \STATE Sample a data point $z_t \sim \mathcal{D}$ uniformly at random
    \STATE Compute stochastic gradient: $g_t \leftarrow \nabla f(w_t; z_t)$
    \STATE Update parameters: $w_{t+1} \leftarrow w_t - \eta_t g_t$
\ENDFOR
\RETURN Output $\bar{w}_T$ chosen uniformly at random from $\{w_0, w_1, \dots, w_{T-1}\}$
\end{algorithmic}
\end{algorithm}

\section*{Dry Run Example}
We apply the algorithm to a deterministic quadratic function to illustrate the mechanics. Note that for a deterministic function, $g_t = \nabla f(w_t)$.
Objective function: $f(w) = (w-3)^2$
Gradient: $\nabla f(w) = 2(w-3)$
Initialization: $w_0 = 0$, Learning rate: $\eta = 0.1$ (constant)

\textbf{Iteration 0 ($t=0$):}
\begin{itemize}
    \item Current parameter: $w_0 = 0$
    \item Gradient computation: $g_0 = 2(0 - 3) = -6$
    \item Update: $w_1 = w_0 - \eta g_0 = 0 - 0.1(-6) = 0.6$
    \item Objective value: $f(w_1) = (0.6 - 3)^2 = (-2.4)^2 = 5.76$
\end{itemize}

\textbf{Iteration 1 ($t=1$):}
\begin{itemize}
    \item Current parameter: $w_1 = 0.6$
    \item Gradient computation: $g_1 = 2(0.6 - 3) = -4.8$
    \item Update: $w_2 = w_1 - \eta g_1 = 0.6 - 0.1(-4.8) = 1.08$
    \item Objective value: $f(w_2) = (1.08 - 3)^2 = (-1.92)^2 = 3.6864$
\end{itemize}

\textbf{Iteration 2 ($t=2$):}
\begin{itemize}
    \item Current parameter: $w_2 = 1.08$
    \item Gradient computation: $g_2 = 2(1.08 - 3) = -3.84$
    \item Update: $w_3 = w_2 - \eta g_2 = 1.08 - 0.1(-3.84) = 1.464$
    \item Objective value: $f(w_3) = (1.464 - 3)^2 = (-1.536)^2 \approx 2.359$
\end{itemize}
The parameters are moving towards the optimal value $w^* = 3$, and the objective value is strictly decreasing.

\section*{Assumptions}
To analyze the convergence of the algorithm for non-convex and smooth functions, we make the following formal assumptions:

\begin{enumerate}
    \item \textbf{Lower Boundedness:} The objective function $f$ is bounded from below. There exists a finite $f^*$ such that $f(w) \geq f^*$ for all $w \in \mathbb{R}^d$.
    \item \textbf{$L$-Smoothness (Lipschitz Continuous Gradient):} The function $f$ is differentiable, and its gradient is $L$-Lipschitz continuous. There exists a constant $L > 0$ such that for all $x, y \in \mathbb{R}^d$:
    $$\|\nabla f(x) - \nabla f(y)\| \leq L\|x - y\|$$
    \item \textbf{Unbiased Stochastic Gradient:} The stochastic gradient is an unbiased estimator of the true gradient. For all $t$:
    $$\mathbb{E}[g_t \mid w_t] = \nabla f(w_t)$$
    \item \textbf{Bounded Variance:} The variance of the stochastic gradient is bounded by a constant $\sigma^2 \geq 0$. For all $t$:
    $$\mathbb{E}[\|g_t - \nabla f(w_t)\|^2 \mid w_t] \leq \sigma^2$$
\end{enumerate}

\section*{Main Convergence Theorem}
\textbf{Theorem 1:} Under Assumptions 1-4, if we run SGD for $T$ iterations with a constant learning rate $\eta_t = \eta \leq \frac{1}{L}$, then the expected squared norm of the gradient of the randomly chosen iterate $\bar{w}_T$ is bounded by:
$$ \mathbb{E}[\|\nabla f(\bar{w}_T)\|^2] \leq \frac{2(f(w_0) - f^*)}{\eta T} + L\eta\sigma^2 $$
Furthermore, by choosing $\eta = \min\left(\frac{1}{L}, \sqrt{\frac{2(f(w_0) - f^*)}{L \sigma^2 T}}\right)$, the algorithm achieves a convergence rate of:
$$ \mathbb{E}[\|\nabla f(\bar{w}_T)\|^2] \leq \mathcal{O}\left(\frac{1}{\sqrt{T}}\right) $$

\section*{Theoretical Properties}
\begin{itemize}
    \item \textbf{Convergence Metric:} Because the function is non-convex, the algorithm is not guaranteed to find a global minimum. Instead, the convergence is measured by the expected squared gradient norm $\mathbb{E}[\|\nabla f(w)\|^2] \rightarrow 0$, which indicates convergence to a first-order stationary point (local minimum, local maximum, or saddle point).
    \item \textbf{Variance Impact:} The term $L\eta\sigma^2$ represents the noise introduced by the stochastic gradients. Unlike deterministic gradient descent which converges to exactly $0$, SGD with a constant step size oscillates around the stationary point with a variance proportional to $\eta$.
    \item \textbf{Hyperparameter Sensitivity:} The learning rate $\eta$ controls the trade-off between the speed of convergence (the $\frac{1}{\eta T}$ term) and the asymptotic noise floor (the $L\eta\sigma^2$ term). A decaying learning rate schedule is necessary for exact convergence to a stationary point.
\end{itemize}

\section*{Discussion}
In the non-convex regime, analyzing SGD diverges significantly from the convex case. Instead of analyzing the distance to the optimum $\|w_t - w^*\|^2$ or the function sub-optimality $f(w_t) - f(w^*)$, we must rely on the Descent Lemma provided by the smoothness assumption. The proof demonstrates that, in expectation, each step of SGD decreases the function value by an amount proportional to the squared gradient magnitude, penalized by the variance of the stochastic approximation.

\section*{Preliminaries (Definitions and Lemmas)}
\textbf{Lemma 1 (Descent Lemma):} For any function $f$ satisfying Assumption 2 ($L$-smoothness), the following inequality holds for any $x, y \in \mathbb{R}^d$:
$$f(y) \leq f(x) + \langle \nabla f(x), y - x \rangle + \frac{L}{2}\|y - x\|^2$$

\textbf{Lemma 2 (Second Moment of Gradient):} Under Assumptions 3 and 4, the expected squared norm of the stochastic gradient is bounded by:
$$\mathbb{E}[\|g_t\|^2 \mid w_t] \leq \|\nabla f(w_t)\|^2 + \sigma^2$$
\textit{Proof:}
$\mathbb{E}[\|g_t\|^2 \mid w_t] = \mathbb{E}[\|g_t - \nabla f(w_t) + \nabla f(w_t)\|^2 \mid w_t] \\
= \mathbb{E}[\|g_t - \nabla f(w_t)\|^2 \mid w_t] + 2\mathbb{E}[\langle g_t - \nabla f(w_t), \nabla f(w_t) \rangle \mid w_t] + \|\nabla f(w_t)\|^2 \\
\leq \sigma^2 + 0 + \|\nabla f(w_t)\|^2$.

\section*{Main Proof (Step-by-Step Derivation)}

We begin by applying Lemma 1 (Descent Lemma) to the iterates $w_t$ and $w_{t+1}$.
\begin{align}
    f(w_{t+1}) &\leq f(w_t) + \langle \nabla f(w_t), w_{t+1} - w_t \rangle + \frac{L}{2} \|w_{t+1} - w_t\|^2 \quad \text{[Step 1]}
\end{align}
Substitute the SGD update rule $w_{t+1} - w_t = -\eta g_t$:
\begin{align}
    f(w_{t+1}) &\leq f(w_t) + \langle \nabla f(w_t), -\eta g_t \rangle + \frac{L}{2} \|-\eta g_t\|^2 \quad \text{[Step 2]} \\
    f(w_{t+1}) &= f(w_t) - \eta \langle \nabla f(w_t), g_t \rangle + \frac{L\eta^2}{2} \|g_t\|^2 \quad \text{[Step 3]}
\end{align}
Take the conditional expectation of both sides with respect to the randomness at step $t$, given the current parameter $w_t$:
\begin{align}
    \mathbb{E}[f(w_{t+1}) \mid w_t] &\leq f(w_t) - \eta \mathbb{E}[\langle \nabla f(w_t), g_t \rangle \mid w_t] + \frac{L\eta^2}{2} \mathbb{E}[\|g_t\|^2 \mid w_t] \quad \text{[Step 4]}
\end{align}
Since $\nabla f(w_t)$ is deterministic given $w_t$, we can pull it out of the inner product expectation:
\begin{align}
    \mathbb{E}[f(w_{t+1}) \mid w_t] &\leq f(w_t) - \eta \langle \nabla f(w_t), \mathbb{E}[g_t \mid w_t] \rangle + \frac{L\eta^2}{2} \mathbb{E}[\|g_t\|^2 \mid w_t] \quad \text{[Step 5]}
\end{align}
Apply Assumption 3 (Unbiasedness, $\mathbb{E}[g_t \mid w_t] = \nabla f(w_t)$):
\begin{align}
    \mathbb{E}[f(w_{t+1}) \mid w_t] &\leq f(w_t) - \eta \langle \nabla f(w_t), \nabla f(w_t) \rangle + \frac{L\eta^2}{2} \mathbb{E}[\|g_t\|^2 \mid w_t] \quad \text{[Step 6]} \\
    \mathbb{E}[f(w_{t+1}) \mid w_t] &\leq f(w_t) - \eta \|\nabla f(w_t)\|^2 + \frac{L\eta^2}{2} \mathbb{E}[\|g_t\|^2 \mid w_t] \quad \text{[Step 7]}
\end{align}
Substitute the bound from Lemma 2 ($\mathbb{E}[\|g_t\|^2 \mid w_t] \leq \|\nabla f(w_t)\|^2 + \sigma^2$):
\begin{align}
    \mathbb{E}[f(w_{t+1}) \mid w_t] &\leq f(w_t) - \eta \|\nabla f(w_t)\|^2 + \frac{L\eta^2}{2} \left( \|\nabla f(w_t)\|^2 + \sigma^2 \right) \quad \text{[Step 8]}
\end{align}
Group the terms involving $\|\nabla f(w_t)\|^2$:
\begin{align}
    \mathbb{E}[f(w_{t+1}) \mid w_t] &\leq f(w_t) - \eta \left( 1 - \frac{L\eta}{2} \right) \|\nabla f(w_t)\|^2 + \frac{L\eta^2\sigma^2}{2} \quad \text{[Step 9]}
\end{align}
Take the total expectation over all randomness up to time $t$:
\begin{align}
    \mathbb{E}[f(w_{t+1})] &\leq \mathbb{E}[f(w_t)] - \eta \left( 1 - \frac{L\eta}{2} \right) \mathbb{E}[\|\nabla f(w_t)\|^2] + \frac{L\eta^2\sigma^2}{2} \quad \text{[Step 10]}
\end{align}
Assume the learning rate is chosen such that $\eta \leq \frac{1}{L}$. This implies $1 - \frac{L\eta}{2} \geq \frac{1}{2}$. Thus:
\begin{align}
    \mathbb{E}[f(w_{t+1})] &\leq \mathbb{E}[f(w_t)] - \frac{\eta}{2} \mathbb{E}[\|\nabla f(w_t)\|^2] + \frac{L\eta^2\sigma^2}{2} \quad \text{[Step 11]}
\end{align}
Rearrange the inequality to isolate the expected squared gradient norm:
\begin{align}
    \frac{\eta}{2} \mathbb{E}[\|\nabla f(w_t)\|^2] &\leq \mathbb{E}[f(w_t)] - \mathbb{E}[f(w_{t+1})] + \frac{L\eta^2\sigma^2}{2} \quad \text{[Step 12]}
\end{align}
Sum the inequality from $t = 0$ to $T - 1$:
\begin{align}
    \sum_{t=0}^{T-1} \frac{\eta}{2} \mathbb{E}[\|\nabla f(w_t)\|^2] &\leq \sum_{t=0}^{T-1} \left( \mathbb{E}[f(w_t)] - \mathbb{E}[f(w_{t+1})] \right) + \sum_{t=0}^{T-1} \frac{L\eta^2\sigma^2}{2} \quad \text{[Step 13]}
\end{align}
The first sum on the right-hand side is a telescoping series:
\begin{align}
    \frac{\eta}{2} \sum_{t=0}^{T-1} \mathbb{E}[\|\nabla f(w_t)\|^2] &\leq f(w_0) - \mathbb{E}[f(w_T)] + \frac{TL\eta^2\sigma^2}{2} \quad \text{[Step 14]}
\end{align}
Using Assumption 1 ($f(w_T) \geq f^*$), we can bound $-\mathbb{E}[f(w_T)] \leq -f^*$:
\begin{align}
    \frac{\eta}{2} \sum_{t=0}^{T-1} \mathbb{E}[\|\nabla f(w_t)\|^2] &\leq f(w_0) - f^* + \frac{TL\eta^2\sigma^2}{2} \quad \text{[Step 15]}
\end{align}
Divide both sides by $\frac{\eta T}{2}$:
\begin{align}
    \frac{1}{T} \sum_{t=0}^{T-1} \mathbb{E}[\|\nabla f(w_t)\|^2] &\leq \frac{2(f(w_0) - f^*)}{\eta T} + L\eta\sigma^2 \quad \text{[Step 16]}
\end{align}
Recall the definition of $\bar{w}_T$ from the pseudocode, which is selected uniformly at random from $\{w_0, \dots, w_{T-1}\}$. The expectation of its squared gradient norm is exactly the average:
\begin{align}
    \mathbb{E}[\|\nabla f(\bar{w}_T)\|^2] &= \frac{1}{T} \sum_{t=0}^{T-1} \mathbb{E}[\|\nabla f(w_t)\|^2] \quad \text{[Step 17]}
\end{align}
Substituting this yields the final bound:
\begin{align}
    \mathbb{E}[\|\nabla f(\bar{w}_T)\|^2] &\leq \frac{2(f(w_0) - f^*)}{\eta T} + L\eta\sigma^2 \quad \text{[Step 18]}
\end{align}

\section*{Rate Derivation (With Constants)}
To optimize the convergence rate, we must select $\eta$ to minimize the right-hand side of the bound from Step 18.
Let $\Delta = f(w_0) - f^*$. We want to minimize:
$$ h(\eta) = \frac{2\Delta}{\eta T} + L\eta\sigma^2 $$
Taking the derivative with respect to $\eta$ and setting it to zero:
$$ h'(\eta) = -\frac{2\Delta}{\eta^2 T} + L\sigma^2 = 0 \implies \eta^* = \sqrt{\frac{2\Delta}{L\sigma^2 T}} $$
However, we must also satisfy the assumption that $\eta \leq \frac{1}{L}$. Thus, we set:
$$ \eta = \min\left(\frac{1}{L}, \sqrt{\frac{2\Delta}{L\sigma^2 T}}\right) $$
\textbf{Case 1: $T$ is small such that $\frac{1}{L} \leq \sqrt{\frac{2\Delta}{L\sigma^2 T}}$} \\
Then $\eta = \frac{1}{L}$. Substituting this into the bound:
$$ \mathbb{E}[\|\nabla f(\bar{w}_T)\|^2] \leq \frac{2\Delta L}{T} + \sigma^2 $$
Here, the algorithm is dominated by the initial optimization phase, and the noise $\sigma^2$ forms a constant floor.

\textbf{Case 2: $T$ is large such that $\sqrt{\frac{2\Delta}{L\sigma^2 T}} \leq \frac{1}{L}$} \\
Then $\eta = \sqrt{\frac{2\Delta}{L\sigma^2 T}}$. Substituting this into the bound:
\begin{align*}
    \mathbb{E}[\|\nabla f(\bar{w}_T)\|^2] &\leq \frac{2\Delta}{T} \sqrt{\frac{L\sigma^2 T}{2\Delta}} + L\sigma^2 \sqrt{\frac{2\Delta}{L\sigma^2 T}} \\
    &= \sqrt{\frac{2\Delta L \sigma^2}{T}} + \sqrt{\frac{2\Delta L \sigma^2}{T}} \\
    &= 2\sqrt{\frac{2\Delta L \sigma^2}{T}}
\end{align*}
This gives the final convergence rate of $\mathcal{O}\left(\frac{1}{\sqrt{T}}\right)$.

\section*{Special Cases}
\textbf{Deterministic Gradient Descent ($\sigma^2 = 0$):}
If full gradients are used, there is no variance ($\sigma^2 = 0$). In this case, the optimal learning rate strictly following the constraint is $\eta = \frac{1}{L}$. The bound from Step 18 simplifies to:
$$ \mathbb{E}[\|\nabla f(\bar{w}_T)\|^2] \leq \frac{2(f(w_0) - f^*) T}{T/L} = \frac{2L(f(w_0) - f^*)}{T} $$
This yields a convergence rate of $\mathcal{O}\left(\frac{1}{T}\right)$, recovering the standard rate for deterministic gradient descent on smooth non-convex functions.

\end{document}